% !PRINTED 25
\weberpage{25}
%<<CONT
denken aus den Functionen einer Veränderlichen $x$, wenn man darin die
Coëfficienten $a_0,a_1\ldots a_n$ selbst wieder als ganze rationale
Functionen von anderen Veränderlichen $y,z\ldots$ auffasst. So
entstehen Functionen von $m$ Veränderlichen aus Functionen von $m-1$
Veränderlichen. Alle Glieder einer solchen Function sind von der Form
$x^r y^s z^t\cdots$, multiplicirt mit einem Coëfficienten, und wenn die
Function gehörig zusammengefasst ist, so kommt jede Combination der
Exponenten $r,s,t\ldots$ nur einmal vor. Zwei so geordnete ganze
Functionen gelten nur dann als einander gleich, wenn sie dieselben
Producte $x^r y^s z^t\cdots$ mit denselben Coëfficienten enthalten.

\begin{center}
{\large \S.\ 2.}

\vspace{0.5em}
{\large\gesperrt{Ein Satz von Gauss.}}
\end{center}

Wir wollen sogleich eine Anwendung der Multiplicationsregel zweier
ganzer rationaler Functionen machen zum Beweise eines Satzes von
Gauss, der uns später noch nützlich sein wird, hier aber zur Einführung
in die Rechnungsweise und als Beispiel dienen soll.\footnote{Gauss,
\emph{Disquisitiones arithmeticae}, Art.~42.}

Wir betrachten hier den Fall, dass die Coëfficienten in den Functionen
$A(x),B(x)$ ganze Zahlen sind, so dass nach (4) auch die Coëfficienten
von $C(x)=A(x)B(x)$ ganze Zahlen sind.

Wenn die sämmtlichen ganzzahligen Coëfficienten
$a_0,a_1,\ldots,a_m$ einer Function $A(x)$ keinen gemeinschaftlichen
Theiler haben, so heisst die Function $A(x)$ eine
\gesperrt{ursprüngliche oder primitive} Function und der Satz, den wir
beweisen wollen, lautet:

\gesperrt{Wenn $A(x)$ und $B(x)$ ursprüngliche Functionen sind, so ist
auch ihr Product $C(x)$ eine ursprüngliche Function.}

Der Beweis ergiebt sich fast unmittelbar aus dem Anblick der Formel
(4).

Wenn nämlich die sämmtlichen Coëfficienten
$c_0,c_1,c_2\ldots c_{m+n}$, wie sie in (4) angegeben sind, einen
gemeinschaftlichen Theiler haben, der grösser als $1$ ist, so muss es
auch wenigstens eine Primzahl geben, die in allen diesen Coëfficienten
aufgeht. Es sei $p$ eine solche Primzahl; diese kann nach der
Voraussetzung, dass $A(x),B(x)$ ursprünglich seien, weder in allen
$a_0,a_1\ldots a_m$, noch in allen $b_0,b_1\ldots b_n$ aufgehen.
